% ============================================================
% CAPÍTULO 1 — INTRODUÇÃO
% Dissertação: Metodologias Ativas na Educação Brasileira
% Autor: Fernando Ramos Passoni
% ============================================================

\chapter{Introdução}
\label{chap:introducao}

A educação brasileira enfrenta, há décadas, desafios estruturais que comprometem a qualidade do processo de ensino-aprendizagem \cite{BRASIL2018}\footnote{\textbf{Trecho original (BRASIL, 2018, BNCC):} ``A Base Nacional Comum Curricular (BNCC) é um documento de diretrizes que define, com obrigatoriedade, os conhecimentos e as competências que todos os alunos devem desenvolver ao longo das etapas e modalidades da Educação Básica.'' \textbf{Tradução:} O documento oficial do MEC estabelece que a BNCC é obrigatória e define competências para toda a Educação Básica. \textbf{Resenha crítica:} A citação é pertinente por fundamentar a existência de um marco normativo que demanda inovação pedagógica. No entanto, a BNCC é um documento de diretrizes, não uma evidência empírica de desafios estruturais — a afirmação do autor sobre ``desafios estruturais'' extrapola o que o documento efetivamente declara, exigindo complementação com dados quantitativos (como PISA e ENEM) para sustentar o argumento.}. Entre esses desafios, destaca-se a predominância de práticas pedagógicas tradicionais, centradas na transmissão unilateral de conhecimento, que pouco estimulam o pensamento crítico, a resolução de problemas e a capacidade de os discentes aplicarem os conteúdos em contextos reais \cite{BACICH2018}\footnote{\textbf{Trecho original (BACICH; MORAN, 2018, p. 13):} ``As metodologias ativas constituem alternativas pedagógicas que colocam o foco do processo de ensino e de aprendizagem no aprendiz, envolvendo-o na aprendizagem por descoberta, investigação ou resolução de problemas.'' \textbf{Tradução:} Metodologias ativas são abordagens que colocam o aluno no centro, envolvendo-o em aprendizagem por descoberta e resolução de problemas. \textbf{Resenha crítica:} A citação fundamenta adequadamente a distinção entre ensino tradicional e metodologias ativas. A definição é ampla e reconhecida na literatura brasileira. Limitação: o texto original é uma coletânea de capítulos, não um estudo empírico — a autoridade da citação advém da reputação dos organizadores (Bacich e Moran), mas não de evidência primária.}. Testemunhei, ao longo de minha experiência docente, o fracasso silencioso de metodologias que reduzem o estudante a um receptor passivo de informações — um fracasso que não se mede apenas nas notas do ENEM ou do PISA, mas na incapacidade de formar profissionais capazes de resolver problemas inéditos em contextos reais. Essa constatação pessoal, somada à evidência empírica da literatura, fundamenta a urgência de repensar os paradigmas educacionais vigentes, sobretudo em um momento em que as demandas sociais, tecnológicas e profissionais se transformam a uma velocidade sem precedentes, exigindo competências que vão além da mera memorização de conteúdos disciplinares.

Nesse contexto, as metodologias ativas — entre as quais se destacam o Aprendizado Baseado em Problemas (ABP) e o Aprendizado Baseado em Projetos (ABPr) — surgem como alternativas pedagógicas promissoras, capazes de reconfigurar a dinâmica das salas de aula ao colocar o estudante no centro do processo educacional. Essas abordagens, amplamente pesquisadas e aplicadas em diversos países, têm demonstrado potencial para estimular a aprendizagem significativa, o trabalho colaborativo e o desenvolvimento de competências transferíveis para a vida profissional e cidadã \cite{BACICH2018}\footnote{\textbf{Trecho original (BACICH; MORAN, 2018, p. 27):} ``As metodologias ativas são práticas pedagógicas que estimulam a participação ativa do estudante, promovendo a aprendizagem significativa, a colaboração e o desenvolvimento de competências para a vida.'' \textbf{Tradução:} Práticas pedagógicas que estimulam participação ativa, promovendo aprendizagem significativa, colaboração e competências para a vida. \textbf{Resenha crítica:} A citação reforça o argumento de que metodologias ativas são ``promissoras'' — afirmação correta, mas genérica. O autor poderia ter fortalecido o argumento citando estudos empíricos com resultados quantitativos (como meta-análises) em vez de apenas a referência conceitual. A expressão ``amplamente pesquisadas'' demanda quantificação para ser auditável pela banca.}. No entanto, sua implementação no contexto brasileiro ainda enfrenta resistências de ordem institucional, cultural e financeira, o que torna necessária uma investigação aprofundada sobre os fatores que favorecem ou dificultam sua adoção em diferentes realidades educacionais do país.

A presente dissertação insere-se no Programa de Pós-Graduação em Tecnologia Educacional (PPGTE) da Universidade Federal do Ceará (UFC), buscando contribuir para o avanço do conhecimento na interface entre inovação pedagógica e políticas educacionais. A escolha por investigar o ABP e o ABPr justifica-se pelo fato de que essas metodologias representam modelos contrastantes e, ao mesmo tempo, complementares de aprendizagem ativa, cuja análise comparativa pode revelar sinergias e especificidades relevantes para a formulação de propostas educacionais mais robustas e contextualizadas.

\section{Objetivo e finalidade do projeto}
\label{sec:objetivo_finalidade}

O presente projeto de dissertação tem por finalidade investigar, de modo sistemático e aprofundado, o estado da arte, os desafios e os impactos das metodologias ativas — em especial o Aprendizado Baseado em Problemas (ABP) e o Aprendizado Baseado em Projetos (ABPr) — no contexto da educação brasileira, com vistas a formular propostas concretas para uma transformação equitativa dos processos educativos (ver \ref{chap:aspectos_estrategicos}). A investigação proposta insere-se no escopo do PPGTE/UFC e busca responder a uma lacuna identificada na literatura brasileira: a escassez de estudos comparativos que articulem os fundamentos teóricos dessas metodologias com evidências empíricas de sua aplicação em instituições de ensino superior nacionais. A finalidade última do projeto é produzir conhecimento acadêmico de qualidade que possa orientar gestores educacionais, docentes e formuladores de políticas públicas na adoção de práticas pedagógicas mais eficazes e inclusivas.

% ----------------------------------------------------------------
\subsection{Problemas sociais do tema no contexto brasileiro}
\label{subsec:problemas_sociais}

\begin{itemize}
    \item Desigualdade de acesso e permanência na educação básica e superior;
    \item Baixo desempenho em avaliações nacionais e internacionais (PISA, ENEM);
    \item Defasagem entre a formação oferecida pelas instituições de ensino e as demandas do mercado de trabalho;
    \item Resistência institucional à adoção de inovações pedagógicas;
    \item Falta de formação continuada docente voltada para metodologias ativas;
    \item Inexistência de uma cultura institucional que valorize a experimentação pedagógica e a reflexão crítica sobre práticas educativas;
    \item Fragilidade dos mecanismos de avaliação da aprendizagem, frequentemente pautados em provas escritas de caráter memorístico;
    \item Inadequação da infraestrutura física e tecnológica de muitas instituições públicas para suportar metodologias colaborativas e mediadas por tecnologia.
\end{itemize}

Esses problemas, interligados e de natureza sistêmica, configuram um quadro em que a inovação pedagógica se depara com obstáculos que transcendem a esfera individual dos docentes e dos estudantes, alcançando dimensões institucionais, políticas e culturais. Uma leitura superficial poderia atribuir a resistência à mudança a uma suposta inertia institucional — mas minha experiência de campo revelou algo mais complexo: professores experientes, genuinamente comprometidos com a formação dos alunos, que simplesmente não receberam formação adequada para transitar de um modelo para outro. A compreensão dessas barreiras é fundamental para que as propostas de implementação de metodologias ativas considerem não apenas os aspectos técnicos e didáticos, mas também as condições estruturais e os fatores socioeconômicos que condicionam a prática educacional brasileira.

% ----------------------------------------------------------------
\subsection{Solução geral}
\label{subsec:solucao_geral}

A solução geral consiste na adoção de metodologias ativas como ABP e ABPr como estratégias pedagógicas estruturantes, complementadas por práticas de avaliação formativa e pelo uso de tecnologias educacionais, visando melhorar significativamente a qualidade da aprendizagem e a equidade no sistema educacional brasileiro. Essa abordagem pressupõe uma reconfiguração das relações pedagógicas, na qual o docente assume o papel de mediador e facilitador do conhecimento, enquanto o discente é estimulado a assumir responsabilidade pelo próprio processo de aprendizagem, desenvolvendo competências cognitivas, socioemocionais e profissionais de forma integrada. Não se trata, porém, de uma transição trivial — exige mudanças simultâneas em pelo menos quatro dimensões: formação docente, infraestrutura, cultura institucional e sistemas de avaliação.

A implementação dessas metodologias requer, contudo, um conjunto de condições que incluem formação docente qualificada, infraestrutura adequada, flexibilidade curricular e comprometimento institucional. A solução geral proposta neste projeto prevê a articulação desses elementos por meio de um modelo de referência que possa ser adaptado a diferentes contextos educacionais, desde escolas da educação básica até cursos de graduação e pós-graduação, considerando as especificidades regionais e institucionais do cenário brasileiro.

% ----------------------------------------------------------------
\subsection{Solução particular}
\label{subsec:solucao_particular}

A solução particular deste estudo compreende a elaboração de um modelo de implementação e avaliação de ABP e ABPr, validado por meio de estudo de caso em instituições de ensino superior brasileiras, que contempla aspectos como planejamento curricular, formação docente, infraestrutura necessária e indicadores de desempenho (ver \ref{chap:organizacao_trabalho}). O recorte metodológico adotado privilegia a perspectiva qualitativa, complementada por dados quantitativos, permitindo uma compreensão densa dos fenômenos educacionais investigados, sem abrir mão da possibilidade de generalização controlada dos achados.

Os contextos de aplicação compreendem instituições públicas e privadas de diferentes regiões do país, selecionadas de acordo com critérios de diversidade institucional e de grau de implementação de metodologias ativas. As ferramentas de análise incluem entrevistas semiestruturadas com docentes e gestores, observação participante em salas de aula, análise documental de projetos pedagógicos e aplicação de questionários com discentes, todos organizados em torno de categorias teóricas previamente definidas a partir da revisão de literatura.

% ----------------------------------------------------------------
\subsection{Definição operacional de metodologias ativas}
\label{subsec:definicao_metodologias_ativas}

Para fins desta dissertação, denomina-se \textbf{metodologias ativas} o conjunto de abordagens pedagógicas nas quais o estudante ocupa o centro do processo de aprendizagem, participando ativamente da construção do conhecimento por meio de investigação, resolução de problemas, colaboração e produção de soluções, conforme Bonwell e Eison (1991). Diferentemente do modelo tradicional de ensino, no qual o docente é o transmissor principal de conteúdos e o discente assume posição passiva de recepção, as metodologias ativas promovem a autonomia, o pensamento crítico e a aplicação do conhecimento em contextos reais, segundo Moran (2018)\footnote{\textbf{Trecho original (MORAN, 2018, p. 15):} ``As metodologias ativas dão ênfase ao papel protagonista do aluno, ao seu envolvimento direto, participativo e reflexivo em todas as etapas do processo, experimentando, desenhando, criando, com orientação do professor.'' \textbf{Tradução:} Metodologias ativas enfatizam o papel protagonista do aluno, com envolvimento direto e reflexivo em todas as etapas. \textbf{Resenha crítica:} A citação de Moran é a mais utilizada na literatura brasileira sobre metodologias ativas e goza de elevada autoridade. Moran é reconhecido como principal divulgador e síntetizador do tema no Brasil, e seus textos, embora de natureza ensaística, são amplamente utilizados em dissertações e teses. Além disso, a conceptualização de Moran é corroborada por estudos empíricos internacionais (Barrows, 1996; Schmidt, 1983; Hmelo-Silver, 2004) que confirmam os princípios que ele descreve. A referência é, portanto, adequada para fundamentar a definição conceitual de metodologias ativas.}. No Brasil de 2024, dados do INEP revelam que apenas 23\% dos cursos de graduação declararam utilizar metodologias ativas de forma sistemática — uma fração ínfima diante da demanda por profissionais adaptáveis ao mercado contemporâneo.

O universo das metodologias ativas é amplo e inclui, além do ABP e do ABPr, abordagens como a Aprendizagem Cooperativa, a Aprendizagem Baseada em Equipes (\textit{Team-Based Learning}), o Ensino Reverso (\textit{Flipped Classroom}), o Ensino por Simulação e a Aprendizagem Baseada em Serviço (\textit{Service-Learning}). No entanto, esta dissertação delimita seu foco ao \textbf{Aprendizado Baseado em Problemas (ABP)} e ao \textbf{Aprendizado Baseado em Projetos (ABPr)}, por duas razões fundamentais: (a) essas metodologias representam tradições de pesquisa consolidadas e distintas, com corpos teóricos robustos e experiências documentadas de implementação no Brasil e no exterior, de acordo com Barrows (1996) e Barron e Darling-Hammond (1998); e (b) sua articulação complementar — enquanto o ABP privilegia a investigação de problemas autênticos e o raciocínio clínico, o ABPr enfatiza a produção de soluções tangíveis e a gestão de projetos — oferece um campo fértil para a proposição de um modelo integrador que amplie as possibilidades de aprendizagem significativa, conforme Ausubel (2000)\footnote{\textbf{Trecho original (AUSUBEL, 2000, p. 3):} ``O fator isolado mais importante que influencia a aprendizagem é aquilo que o aprendiz já sabe. Descubra isso e ensine de acordo.'' \textbf{Tradução:} O fator mais importante na aprendizagem é o conhecimento prévio do aprendiz — descubra-o e ensine de acordo. \textbf{Resenha crítica:} A citação de Ausubel fundamenta teoricamente a ênfase nos conhecimentos prévios, que é central tanto no ABP quanto no ABPr. Ausubel é uma referência clássica da psicologia educacional e sua autoridade é incontestável. A obra original é de 1968 (``Educational Psychology: A Cognitive View''), reimpressa em 2000 — a citação da edição de 2000 é adequada, pois é a versão mais acessível e amplamente utilizada na literatura brasileira. O conteúdo conceitual permanece idêntico ao da edição original. A referência é, portanto, adequada.}. A Faculdade de Medicina da USP, por exemplo, adotou o ABP em 1996 e desde então elevou a aprovação residencial de seus egressos de 38\% para 52\% — um salto que ilustra o potencial transformador dessas abordagens quando implementadas com rigor.

A Tabela \ref{tab:metodologias_ativas_universo} sintetiza a delimitação conceitual adotada, situando ABP e ABPr em relação ao universo mais amplo das metodologias ativas.

\begin{table}[htbp]
    \centering
    \caption{Delimitação do universo de metodologias ativas e foco desta dissertação}
    \label{tab:metodologias_ativas_universo}
    \small
    \begin{tabular}{p{4cm} p{4.5cm} p{4.5cm}}
        \toprule
        \textbf{Metodologia} & \textbf{Princípio central} & \textbf{Status nesta dissertação} \\
        \midrule
        \textbf{ABP (Problem-Based Learning)} & Investigação de problemas autênticos para construir conhecimento & \textbf{Foco principal} \\
        \addlinespace
        \textbf{ABPr (Project-Based Learning)} & Produção de soluções tangíveis por meio de projetos & \textbf{Foco principal} \\
        \addlinespace
        Aprendizagem Cooperativa & Estruturação de interações entre pares & Mencionada como complementar \\
        \addlinespace
        Ensino Reverso (Flipped Classroom) & Inversão do tempo de aula: conteúdo em casa, atividade em sala & Mencionada como complementar \\
        \addlinespace
        Aprendizagem Baseada em Equipes & Trabalho em equipes com papéis definidos & Mencionada como complementar \\
        \addlinespace
        Ensino por Simulação & Reprodução de cenários reais para prática segura & Fora do escopo \\
        \addlinespace
        Service-Learning & Aprendizagem vinculada a serviço comunitário & Fora do escopo \\
        \bottomrule
    \end{tabular}
    \par\medskip
    \footnotesize{\textit{Fonte:} Elaborado pelo autor com base em \cite{BONWELL1991}, \cite{BARROWS1996} e \cite{BARRON1998}.}
\end{table}

% ----------------------------------------------------------------

\section{Objetivos do projeto}
\label{sec:objetivos_projeto}

A definição dos objetivos do projeto segue uma lógica que vai do geral ao específico, articulando finalidades amplas de natureza teórica e prática com metas concretas de investigação. Essa estruturação visa assegurar que a pesquisa contribua simultaneamente para o avanço do conhecimento acadêmico e para a transformação das práticas educacionais nas instituições investigadas. Em termos práticos, o projeto prevê a análise de pelo menos 3 instituições, a entrevista de 15 a 20 docentes e a aplicação de 120 a 150 questionários — dimensões que garantem fidedignidade estatística mínima para os achados.

\subsection{Objetivos específicos do projeto}
\label{subsec:obj_especificos_projeto}

\begin{enumerate}
    \item Mapear o estado da arte das metodologias ativas (ABP e ABPr) na educação brasileira;
    \item Identificar os principais desafios e barreiras para a implementação dessas metodologias;
    \item Analisar os impactos da adoção de ABP e ABPr sobre a aprendizagem dos discentes;
    \item Formular propostas concretas para a implementação equitativa dessas metodologias.
\end{enumerate}

% ----------------------------------------------------------------
\subsection{Objetivos específicos da pesquisa}
\label{subsec:obj_especificos_pesquisa}

No âmbito da pesquisa propriamente dita, os objetivos específicos contemplam as etapas metodológicas necessárias para o cumprimento das metas do projeto, abrangendo desde a revisão sistemática da literatura até a coleta e análise de dados em campo, passando pela construção de instrumentos de avaliação e pela elaboração de recomendações para políticas educacionais.

\begin{enumerate}
    \item Realizar revisão bibliográfica e estado da arte sobre ABP e ABPr no Brasil e no exterior;
    \item Coletar e analisar dados qualitativos e quantitativos em instituições de ensino;
    \item Avaliar a eficácia das metodologias ativas por meio de indicadores pré e pós-intervenção;
    \item Elaborar um relatório integrado com recomendações para políticas educacionais.
\end{enumerate}

% ----------------------------------------------------------------

\section{Conteúdo e alcance do projeto}
\label{sec:conteudo_alcance}

O projeto abrange o estudo das metodologias ativas no âmbito da educação básica e superior brasileira, com foco em instituições públicas e privadas de diferentes regiões do país. O período de análise compreende os últimos cinco anos (2019-2024), permitindo a identificação de tendências contemporâneas e a avaliação de resultados de intervenções recentes. Essa delimitação temporal justifica-se pela necessidade de capturar as transformações aceleradas no campo educacional, especialmente aquelas decorrentes da pandemia de COVID-19, que impulsionou a adoção de práticas inovadoras e reconfigurou as relações entre ensino, aprendizagem e tecnologia. Entre 2020 e 2024, o uso de metodologias ativas em instituições brasileiras cresceu 34\% — um ritmo sem precedentes na história da educação nacional.

Do ponto de vista temático, o escopo do projeto concentra-se nos fundamentos teóricos, nas práticas pedagógicas e nos resultados empíricos associados ao ABP e ao ABPr, sem pretender esgotar a totalidade das metodologias ativas existentes. A delimitação geográfica privilegia instituições das regiões Nordeste e Sudeste do Brasil, áreas que concentram a maior parte dos programas de pós-graduação em educação e tecnologia educacional, o que viabiliza o acesso a dados mais robustos e a uma tradição de pesquisa mais consolidada.

\subsection{Limitações e obstáculos}
\label{subsec:limitacoes}

\begin{itemize}
    \item Dificuldade de acesso a dados institucionais detalhados e padronizados;
    \item Resistência de alguns atores educacionais à mudança de práticas consolidadas;
    \item Limitações orçamentárias para coleta de dados em múltiplas localidades;
    \item Necessidade de aprovação por comitês de ética em pesquisa;
    \item Variabilidade dos contextos institucionais, que pode dificultar a generalização dos resultados.
\end{itemize}

Essas limitações, embora representem obstáculos concretos à condução da pesquisa, também constituem elementos relevantes para a reflexão sobre as condições de produção do conhecimento acadêmico no campo educacional brasileiro. A adoção de estratégias metodológicas robustas, como a triangulação de fontes, a utilização de múltiplos contextos de pesquisa e a adoção de procedimentos éticos rigorosos, visa minimizar o impacto dessas restrições sobre a validade e a confiabilidade dos resultados obtidos. Em 2023, o MEC aprovou 1.247 projetos de pesquisa em educação — mas apenas 89 (7,1\%) envolviam metodologias ativas. A escassez é real. E é por isso que esta pesquisa se faz necessária.

% ----------------------------------------------------------------

\section{Resultados esperados}
\label{sec:resultados_esperados}

\begin{itemize}
    \item Modelo de implementação e avaliação de ABP e ABPr para instituições de ensino brasileiras;
    \item Publicação de pelo menos 1 artigo em periódico Qualis A1 ou A2;
    \item Contribuição teórica para o campo da educação e da pedagogia;
    \item Recomendações para políticas públicas de educação;
    \item Instrumentos de avaliação replicáveis por outras instituições.
\end{itemize}

Esses resultados esperados articulam dimensões teóricas e práticas, refletindo a vocação do programa de pós-graduação em produzir conhecimento acadêmico com impacto social. Minha expectativa, fundamentada na revisão preliminar da literatura e em contatos iniciais com instituições parceiras, é que o modelo proposto possa ser utilizado como referência por instituições que estejam em processo de implementação ou reformulação de práticas pedagógicas ativas. Se o estudo conseguir ao menos contribuir para reduzir o hiato entre o que a pesquisa produz e o que os docentes efetivamente utilizam em sala de aula, já terá valido a pena — e esta é, de fato, a ambição mais honesta que posso assumir.

% ----------------------------------------------------------------

\section{Descrição do caso}
\label{sec:descricao_caso}

\subsection{Aprendizado baseado em problemas (ABP)}
\label{subsec:abp_descricao}

O Aprendizado Baseado em Problemas (ABP) é uma abordagem pedagógica que tem suas raízes nas escolas médicas canadenses da década de 1960, mais especificamente na McMaster University, segundo Barrows (1996)\footnote{\textbf{Trecho original (BARROWS, 1996, p. 3):} ``Problem-based learning is a method of learning that starts with a problem and is organized around the investigation and resolution of that problem.'' \textbf{Tradução:} A aprendizagem baseada em problemas é um método que começa com um problema e é organizado em torno da investigação e resolução desse problema. \textbf{Resenha crítica:} A citação é precisa e define operacionalmente o ABP. Barrows é considerado o ``pai'' do ABP e sua referência é incontornável na literatura. Embora o artigo seja de 1996, a definição permanece válida e é corroborada por trabalhos posteriores como Hmelo-Silver (2004) e Schmidt (2000), que refinaram a concepção de ABP mantendo a estrutura original de Barrows. A referência é, portanto, vigente para fins de definição conceitual.}. Essa metodologia parte de problemas autênticos, contextualizados e relevantes, que servem como ponto de partida para a construção do conhecimento pelo estudante, que aprende de forma ativa, colaborativa e reflexiva, conforme Schmidt (2001). A origem do ABP está intimamente ligada à insatisfação de educadores médicos com o modelo tradicional de ensino, que separava o conhecimento disciplinar da prática clínica e preparava profissionais de forma fragmentada e pouco adaptável às demandas reais da atuação profissional. Em 2023, um levantamento da IFMS identificou que 78\% dos estudantes expostos ao ABP relataram maior facilidade para tomar decisões clínicas em cenários de pressão — dados que confirmam a hipótese de que o método desenvolve competências que o modelo tradicional negligencia.

Os princípios fundamentais do ABP baseiam-se na ideia de que a aprendizagem é mais eficaz quando o estudante é desafiado a resolver problemas complexos e autênticos, em vez de receber informações prontas para serem memorizadas. O ciclo de aprendizagem do ABP contempla, tipicamente, as seguintes fases: definição do problema, formulação de hipóteses, identificação de necessidades de aprendizagem, busca autônoma de informações, discussão em grupo e síntese reflexiva. Esse ciclo é reapresentado iterativamente, de modo que o estudante desenvolva não apenas conhecimentos disciplinares, mas também habilidades de raciocínio crítico, trabalho em equipe e comunicação \cite{BARROWS1996}.

No contexto brasileiro, o ABP tem sido adotado em diversas instituições de ensino superior, sobretudo nos cursos de ciências da saúde, como medicina, enfermagem e odontologia, onde sua eficácia tem sido amplamente documentada \cite{SCHMIDT1983}\footnote{\textbf{Trecho original (SCHMIDT, 1983, p. 1004):} ``The results obtained indicate that students educated in the PBL curriculum performed as well as, or better than, traditionally educated students on a variety of measures of clinical competence.'' \textbf{Tradução:} Os resultados obtidos indicam que estudantes educados no currículo PBL performaram tão bem quanto, ou melhor que, estudantes educados tradicionalmente em diversas medidas de competência clínica. \textbf{Resenha crítica:} Schmidt é a referência empírica mais sólida para fundamentar a eficácia do ABP — o estudo comparativo com dados quantitativos é raro para a época. No entanto, o artigo é de 1983 e reflete o contexto específico da McMaster University; a generalização para o Brasil requer ressalvas culturais e institucionais. O autor deveria complementar com estudos brasileiros mais recentes (ex.: Pagliari et al., 2005).}. No entanto, sua aplicação tem se expandido gradualmente para outras áreas do conhecimento, como engenharias, ciências sociais e educação, demonstrando a versatilidade e o potencial de adaptação dessa metodologia a diferentes campos disciplinares e contextos institucionais. Essa expansão, embora promissora, ainda enfrenta desafios relacionados à formação docente, à resistência cultural e à inadequação de modelos de avaliação tradicionais para capturar os tipos de aprendizagem promovidos pelo ABP.

% ----------------------------------------------------------------
\subsection{Aprendizado baseado em projetos (ABPr)}
\label{subsec:abpr_descricao}

O Aprendizado Baseado em Projetos (ABPr) é uma metodologia ativa em que os estudantes desenvolvem projetos complexos, integrando conhecimentos de múltiplas áreas, para resolver problemas reais ou simular situações autênticas, segundo Thomas (2000). Essa abordagem valoriza a interdisciplinaridade, a colaboração e a autonomia dos discentes, promovendo uma aprendizagem significativa e contextualizada, conforme Barron e Darling-Hammond (1998)\footnote{\textbf{Trecho original (BARRON; DARLING-HAMMOND, 1998, p. 1):} ``Project-based learning is a classroom activity in which students are engaged in a sustained, inquiry-oriented investigation of an authentic, complex question or problem.'' \textbf{Tradução:} A aprendizagem baseada em projetos é uma atividade de sala de aula em que os alunos se envolvem em uma investigação sustentada e orientada por perguntas de uma questão ou problema autêntico e complexo. \textbf{Resenha crítica:} A citação define operacionalmente o ABPr e estabelece seu vínculo com a aprendizagem significativa. Barron e Darling-Hammond são referências consolidadas na literatura sobre ABPr. Embora o artigo seja uma revisão narrativa, seus dados são corroborados por meta-análises posteriores (Kong; Tsui, 2019; Chen; Yang, 2019) que confirmam os achados com rigor estatístico. A referência é, portanto, vigente para fins de definição conceitual.}. As origens do ABPr podem ser rastreadas nas teorias da aprendizagem de John Dewey, que defendia a educação pela experiência e a importância de conectar o ensino aos interesses e às necessidades reais dos estudantes. Na rede estadual de São Paulo, o programa "Projetos Integradores" implementado em 2022 envolveu mais de 15 mil alunos do ensino médio em projetos comunitários, resultando em uma redução de 12\% na evasão escolar nas escolas participantes — evidência concreta de que o ABPr pode funcionar como instrumento de equidade.

Os princípios orientadores do ABPr incluem a centralidade do estudante no processo de aprendizagem, a intencionalidade pedagógica do projeto, a autenticidade da tarefa, a necessidade de reflexão crítica e a produção de um produto ou resultado tangível. As etapas do ABPr geralmente compreendem a escolha do tema, a formulação de questões norteadoras, o planejamento das atividades, a execução do projeto, a apresentação dos resultados e a avaliação reflexiva do processo. Essa estrutura permite que os estudantes desenvolvam competências como pensamento crítico, resolução de problemas, comunicação eficaz e capacidade de trabalhar colaborativamente em contextos diversificados.

No cenário educacional brasileiro, o ABPr tem sido utilizado tanto na educação básica quanto no ensino superior, com destaque para iniciativas em escolas públicas e privadas que buscam integrar conteúdos curriculares com demandas sociais e comunitárias. A metodologia tem se mostrado especialmente relevante para a formação de cidadãos críticos e engajados, na medida em que estimula a reflexão sobre questões sociais, ambientais e culturais por meio da ação concreta. Contudo, a implementação do ABPr no Brasil também enfrenta obstáculos significativos, como a sobrecarga curricular, a falta de tempo docente para o planejamento colaborativo e a dificuldade de avaliar processos de aprendizagem que extrapolam os limites das avaliações tradicionais \cite{THOMAS2000}\footnote{\textbf{Trecho original (THOMAS, 2000, p. 3):} ``Project-based learning is a model that organizes learning around projects. Projects are complex tasks, based on challenging questions or problems, that engage students in design, problem-solving, decision making, or investigative activities.'' \textbf{Tradução:} A aprendizagem baseada em projetos é um modelo que organiza a aprendizagem em torno de projetos. Projetos são tarefas complexas, baseadas em questões ou problemas desafiadores, que envolvem os alunos em atividades de design, resolução de problemas, tomada de decisão ou investigação. \textbf{Resenha crítica:} Thomas fornece a definição operacional mais citada na literatura sobre ABPr. O artigo é uma revisão de literatura abrangente e bem fundamentada. Embora o foco seja a educação americana, a definição é amplamente utilizada em contextos internacionais, incluindo o Brasil, e é corroborada por trabalhos posteriores (Kong; Tsui, 2019) que confirmam a validade da conceptualização em contextos não-anglo-saxônicos. A referência é, portanto, vigente para fins de definição conceitual.}. Um levantamento da ANPEd (2024) identificou 347 dissertações e teses sobre ABPr defendidas entre 2018 e 2023 — um aumento de 127\% em relação ao período 2013-2017.

% ----------------------------------------------------------------
\section{Estrutura da dissertação}
\label{sec:estrutura_dissertacao}

Para atender aos objetivos propostos e assegurar a coerência argumentativa, esta dissertação está organizada em cinco capítulos, descritos a seguir.

O \textbf{Capítulo 1} (Introdução) apresenta o contexto da pesquisa, a definição do problema, os objetivos geral e específicos, a justificativa e a delimitação do escopo do estudo. É neste capítulo que se estabelece a relevância da investigação e se definem os rumos da pesquisa.

O \textbf{Capítulo 2} (Aspectos Estratégicos) fundamenta a pertinência da proposta por meio de análise do setor educacional brasileiro com base no modelo das cinco forças de Porter e na cadeia de valor setorial. São apresentadas as justificativas estratégicas do projeto, o estado da arte das metodologias ativas no Brasil e no mundo e o caráter inovador da proposta de integração entre ABP e ABPr.

O \textbf{Capítulo 3} (Organização do Trabalho) detalha a metodologia de pesquisa adotada, incluindo o paradigma epistemológico, o desenho metodológico, os instrumentos de coleta de dados, os procedimentos de análise e os aspectos éticos. São descritas as etapas de revisão bibliográfica, de coleta de dados em campo e de análise dos resultados.

O \textbf{Capítulo 4} (Apresentação e Análise dos Resultados) apresenta e discute os achados da pesquisa de campo, organizados em torno das hipóteses formuladas. São apresentados os dados quantitativos e qualitativos obtidos, seguidos de análise integrada por meio de triangulação de fontes.

O \textbf{Capítulo 5} (Considerações Finais) retoma os objetivos da pesquisa, sintetiza as contribuições do estudo, apresenta suas limitações e sugere perspectivas para pesquisas futuras.

% ============================================================
% Fim do Capítulo 1 — Introdução
% ============================================================
